\section{Votum}
\label{sec:votum}

Schwarz sind in diesem Teil nur Zahlen und Verweise auf die Teile~\ref{sec:unternehmen}
bis~\ref{sec:flow}; neue Zahlen kommen nicht hinzu. \bk{Blau ist die Abw\"agung,} \vd{rot
das Votum: eine Empfehlung unter den erkl\"arten Annahmen, kein Befund.}

\subsection{Verdikt}

\vd{\textbf{Nicht kaufen zu \je{\Kurs}{EUR}, auf keinem der drei Horizonte.} Auf zw\"olf
Monate liegt der Kurs bereits auf dem Konsensziel, auf drei Jahre \"uber dem Wert, den das
eigene Median-Multiple nach dem Ausbau in Rotterdam ergibt, und auf f\"unf bis f\"unfzehn
Jahre \"uber jedem Schwellenkurs dieses Berichts.}

\begin{table}[htbp]\centering\footnotesize
\caption{Das Votum nach Horizont. Werte je Aktie in EUR, Abstand zum Kurs vom \KursDatum.
Herleitung: kurzer und mittlerer Horizont Tabelle~\ref{tab:kurzfrist}, langer Horizont
Tabelle~\ref{tab:schwelle}.}\label{tab:horizont}
\begin{tabularx}{\linewidth}{@{}p{2.3cm}Xrl@{}}\toprule
Horizont & Ma{\ss}stab & Wert & Lesart\\\midrule
bis 12 Monate & Konsensziel (Sekund\"arquelle) & \KonsensZiel & kein Aufschlag\\
 & Median-Multiple auf EBITDA der letzten zw\"olf Monate & \KurzfristLTM & Marge h\"alt
 (\KurzfristAbstandLTM)\\
 & Median-Multiple auf Median-EBITDA 2019--2025 & \KurzfristMED & Marge normalisiert
 (\KurzfristAbstandMED)\\\addlinespace
1 bis 3 Jahre & Median-Multiple auf Median-EBITDA, skaliert auf die Kapazit\"at nach Rotterdam
 & \KurzfristROT & \KurzfristAbstandROT\\\addlinespace
5 Jahre & Schwellenkurs, LTM ohne Umlaufverm\"ogen / Basis & \SchwelleLTMOFuenf{} / \SchwelleMEDFuenf
 & darunter kaufen\\
10 Jahre & dasselbe; bei \Pz{\HuerdeLang}: \SchwelleLTMOZehnLang & \SchwelleLTMOZehn{} / \SchwelleMEDZehn
 & darunter kaufen\\
15 Jahre & dasselbe & \SchwelleLTMOFuenfzehn{} / \SchwelleMEDFuenfzehn & darunter kaufen\\
\bottomrule\end{tabularx}\end{table}

\vd{\textbf{F\"ur den Halter: reduzieren.} Der Kurs hat seit Jahresende \KursSeitJahresende{}
gewonnen und steht auf dem mittleren Ziel der Analysten; er ruht auf einer Marge am oberen
Ende ihrer eigenen Reihe, die das Unternehmen selbst einem au{\ss}ergew\"ohnlichen Umfeld
zuschreibt. Wer einen Teil verkauft, sichert den Zyklusgewinn; wer den Rest h\"alt, wettet
bewusst auf eine Sondermarge \"uber 2026 hinaus und sollte den Bericht zum dritten Quartal
am 29.~Oktober 2026\Q{GBXXV}{238} als Pr\"ufpunkt setzen.}

\vd{\textbf{F\"ur den Nichthalter: nicht kaufen, Kaufmarken setzen.} F\"ur einen Anleger mit
Horizont bis drei Jahre liegt die Marke bei \je{\KurzfristMED}{EUR}, dem Wert bei
normalisierter Marge und gewohntem Multiple; f\"ur einen Anleger mit zehn Jahren Horizont bei
\je{\SchwelleLTMOZehn}{EUR}, bei der milderen H\"urde \je{\SchwelleLTMOZehnLang}{EUR}.}

\subsection{Begr\"undung aus den Befunden}

\textbf{Was dagegen spricht.} \bk{Die Verkaufsmarge des ersten Halbjahrs 2026 liegt \"uber
jedem Jahr der eigenen Reihe (Teil~\ref{sec:abschluss}), und der Vorstandsvorsitzende
f\"uhrt sie auf den Konflikt im Nahen Osten zur\"uck (Teil~\ref{sec:strategie}). Selbst auf
diesem Niveau erreicht der freie Zufluss nur \FcfRenditeLtm{} des B\"orsenwerts, ohne die
Bewegung des Umlaufverm\"ogens \FcfRenditeLtmOhne. Bei einem Ausstieg zum heutigen
B\"orsenwert verlangt die H\"urde einen Zufluss von \Mio{\ZuflussBedarf}{EUR} j\"ahrlich,
bei der milderen H\"urde \Mio{\ZuflussBedarfLang}{EUR}; das beste Jahr der Reihe brachte
\Mio{\FcfBestesJahr}{EUR} (\FcfBestesJahrJahr). Keine Dauer der Sondermarge schlie{\ss}t diese
L\"ucke (Tabelle~\ref{tab:dauer}). Eine Preissetzungsmacht zeigen die Zahlen nicht: In den
schwachen Jahren lag die eigene Marge unter der Marktreferenz (Tabelle~\ref{tab:wettbewerb}).}

\textbf{Was daf\"ur spricht.} \bk{Rotterdam erh\"oht die Kapazit\"at um \KapazitaetPlusPz{},
und danach sollen die Investitionen auf rund \je{\InvestNachRotterdam}{Mrd.\,EUR} j\"ahrlich
fallen (Teil~\ref{sec:strategie}). Das Leistungsprogramm hat sein Ziel von
\Mio{\ProgrammZiel}{EUR} mit \Mio{\ProgrammLaufrateHj}{EUR} vorzeitig \"ubertroffen, der
Verschuldungsgrad liegt mit \Pz{\VerschuldungHj} unter dem Ziel, und die Nachfrage wird durch
Beimischungspflichten gest\"utzt. Auf das EBITDA der letzten zw\"olf Monate ist die Aktie mit
\EvHeuteLtm{}$\times$ g\"unstiger bewertet als im Median ihrer Reihe (\EvMedian$\times$).}

\textbf{Was f\"ur Warten spricht.} \bk{Die offene Gr\"o{\ss}e ist die H\"ohe der Marge nach
dem Konflikt, und Information dar\"uber kommt schnell: drei Stillst\"ande im zweiten Halbjahr,
der Bericht zum dritten Quartal, die Investitionsprognose f\"ur 2027. Option~B
(Teil~\ref{sec:optionen}) kostet bis dahin nichts.}

\textbf{Wo dieses Votum vom Konsens abweicht und warum.} \bk{In der Richtung weicht es ab:
Der Konsens lautet Kaufen. In der H\"ohe weniger, als es scheint: Das mittlere Kursziel von
\je{\KonsensZiel}{EUR} liegt auf dem heutigen Kurs, die Analysten sehen also selbst keinen
nennenswerten Aufschlag, und ihr niedrigstes Ziel (\je{\KonsensTief}{EUR}) liegt noch unter dem
Normalisierungswert dieses Berichts (\je{\KurzfristMED}{EUR}). Dass die gesamte Abdeckung \"uber den langfristigen
Schwellenkursen liegt, erkl\"art der Mechanismus, nicht ein Irrtum von \KonsensAnalysten{}
H\"ausern: Ein Zw\"olfmonatsziel bewertet das Ergebnis des laufenden und des n\"achsten Jahres
mit einem Multiple und unterstellt einen Fortbestand zum Marktwert; dieser Bericht rechnet mit
dem aussch\"uttbaren Zufluss und dem Buchwert als Endwert. Die Differenz ist die
Bewertungspr\"amie, nicht ein Unterschied in den Zahlen. Sie geht als Gegenargument in
Abschnitt~\ref{sub:grenzen} ein.}

\subsection{Einstiegsschwelle}

\schwellentabelle
\kurzfristtabelle

Die Tabelle~\ref{tab:schwelle} beantwortet die Frage f\"ur den langen Horizont, die
Tabelle~\ref{tab:kurzfrist} f\"ur den kurzen. \bk{Welche Zeile die Daten st\"utzen, ist der
eigentliche Befund: F\"ur den langen Horizont st\"utzt die eigene Reihe den Basisfall
(\je{\SchwelleMEDZehn}{EUR}); das LTM-Szenario ohne Umlaufverm\"ogen ist der g\"unstigste Fall,
den die Berichte hergeben, und liegt mit \je{\SchwelleLTMOZehn}{EUR} noch immer weit unter dem
Kurs. F\"ur den kurzen Horizont stehen die Zeilen f\"ur zwei Welten: Bleibt die Marge, sind
\je{\KurzfristLTM}{EUR} m\"oglich (\KurzfristAbstandLTM), normalisiert sie sich,
\je{\KurzfristMED}{EUR} (\KurzfristAbstandMED). Das Chancen-Risiko-Verh\"altnis ist
asymmetrisch zulasten des K\"aufers.}

\annahme{Die Zeile \enquote{Median, skaliert} unterstellt, dass das vergleichbare EBITDA
proportional zur Kapazit\"at w\"achst, bei Median-Marge und Median-Multiple. Das ist keine
Prognose; es ist die g\"unstigste Lesart des Ausbaus, ohne eine eigene Marge zu erfinden.}

\textbf{Das Multiple im Kontext.}

\evtabelle

\begin{table}[htbp]\centering\small
\caption{Kurs-Gewinn-Verh\"altnis und Vergleich mit zwei Wettbewerbern. Neste-Reihe aus
Tabelle~\ref{tab:ergebnis} (Jahresende, nur positive Ergebnisse, \KgvJahre{} Jahre);
Vergleichswerte Sekund\"arquelle (stockanalysis.com, 27.09.2026), W\"ahrungen je Titel.}\label{tab:multiple}
\begin{tabular}{@{}lrrrr@{}}\toprule
 & KGV nachlaufend & KGV erwartet & EV/EBITDA & KBV\\\midrule
Neste & \PeerNesteKgv & \PeerNesteKgvFwd & \PeerNesteEv & \PeerNesteKbv\\
Valero (50\,\% an Diamond Green Diesel) & \PeerVloKgv & \PeerVloKgvFwd & \PeerVloEv & \PeerVloKbv\\
Darling Ingredients (50\,\%) & \PeerDarKgv & \PeerDarKgvFwd & \PeerDarEv & \PeerDarKbv\\\midrule
Neste, Median der eigenen Reihe & \KgvMedian & -- & \EvMedian & --\\
Neste, auf Konsens 2026 / 2027 & -- & \KgvXXVI{} / \KgvXXVII & -- & --\\
\bottomrule\end{tabular}\end{table}

Das nachlaufende KGV von \KgvLtm{} liegt auf Rang \Pz{\KgvRangLtm} der eigenen Reihe, das
Verh\"altnis von Unternehmenswert zu EBITDA auf Rang \Pz{\EvRangLtm}. Auf das Konsensergebnis
2027 steigt das KGV von \KgvXXVI{} auf \KgvXXVII. Beide Wettbewerber sind auf erwartete
Ergebnisse niedriger bewertet. \bk{Das Multiple steht im unteren Teil seiner Reihe, die
Ertragsgr\"o{\ss}e darunter im oberen: Die beiden Perzentile zeigen in entgegengesetzte
Richtungen. Das ist die bekannte Falle des g\"unstigen Multiples auf einem H\"ochstergebnis.}

\subsection{Ausl\"oser, die das Votum drehen}

\begin{description}\small
\item[Kurs unter \je{\KurzfristMED}{EUR}.] F\"ur Anleger mit Horizont bis drei Jahre neu
  rechnen; unter \je{\SchwelleLTMOZehnLang}{EUR} auch f\"ur den langen Horizont.
\item[Marge im dritten Quartal 2026.] Bericht am 29.~Oktober 2026. H\"alt die Verkaufsmarge
  von Renewable Products trotz der Stillst\"ande auf dem Niveau des ersten Halbjahrs
  (\je{\HjRpMarge}{USD/t}), verschiebt sich die Lesart vom Ausnahmejahr zum neuen Niveau;
  f\"allt sie in Richtung des Medians (\je{\MargeMedian}{USD/t}), ist die Normalisierung
  best\"atigt.
\item[Investitionsprognose 2027.] Liegt sie nahe \je{\InvestNachRotterdam}{Mrd.\,EUR}, beginnt
  die Erntephase planm\"a{\ss}ig; eine erneute Verschiebung oder Verteuerung Rotterdams
  verl\"angert die Bauphase.
\item[Konsensergebnis 2027.] Steigt es \"uber das Ergebnis 2026 (\je{\KonsensEpsXXVI}{EUR}),
  erwartet auch der Markt keinen R\"uckgang mehr; dann tr\"agt das g\"unstige Multiple.
\item[Verschuldungsgrad \"uber \Pz{\Verschuldungsziel}.] Wiederholung von 2024; Votum
  bleibt, Kaufmarken sinken.
\end{description}

\subsection{Was das Votum nicht wissen kann}
\label{sub:grenzen}

\bk{Die Gr\"o{\ss}e, an der alles h\"angt, ist die Marge von Renewable Products nach dem
Konflikt. Sie h\"angt an Politik (Beimischungspflichten in der EU und den USA, Steuergutschriften,
Zertifikatspreise) und an einem geopolitischen Ereignis; beides ist nicht berechenbar.}

\bk{Gegen das eigene Votum sprechen drei Argumente. Erstens der Endwert: Der Buchwert ist
konservativ f\"ur ein Unternehmen, das zu dem \KBV-fachen notiert; wer einen Ausstieg zum
heutigen B\"orsenwert annimmt, erreicht mit der Sondermarge immerhin \DauerIrrMarktZehn{}
p.\,a.\ (Tabelle~\ref{tab:dauer}), noch immer unter beiden H\"urden, aber weit weg von den
negativen Werten der Buchwertrechnung. Zweitens das Wachstum: Die Rechnung h\"alt den Zufluss
konstant, Rotterdam und die Pflichtquoten sind aber Wachstum mit Datum; die Umkehrung zeigt,
wie viel davon der Kurs schon enth\"alt. Drittens der Markt: \KonsensAnalysten{} H\"auser raten
zum Kauf, mehrere haben ihre Ziele nach dem Halbjahr angehoben, und die Aktie hat vom Tief das
\KursTiefZuHeute-fache erreicht. Ein Teil dessen, was dieses Votum abwartet, kann eintreten,
w\"ahrend es wartet.}
